\documentclass[a4paper,12pt]{article}
\usepackage[margin=1.5cm]{geometry}
\usepackage{amsmath, amssymb}
\usepackage{array}
\usepackage{alphalph}

\usepackage{amsmath, amssymb, amsfonts}
\usepackage{multicol}
\usepackage{multirow}

\usepackage{enumitem}
\usepackage{booktabs}
\usepackage{tabularx}
\usepackage{array}
\usepackage{framed}
\usepackage{xcolor}
\usepackage{tikz}
\usepackage{fancyhdr}

% Estilos para cajas de teoría y recordatorios
\definecolor{graybox}{gray}{0.92}
\newenvironment{boxsummary}{%
  \def\FrameCommand{\fcolorbox{black}{graybox}}%
  \MakeFramed{\FrameRestore}%
  \small\textbf{RECORDAR:}\par\smallskip
}{%
  \endMakeFramed
}

\begin{document}

\section*{Tema 2: Potencias y raíces}

\begin{boxsummary}
\[ a^n = a \cdot a \cdot a \cdot \dots \cdot a \cdot a \quad \text{(n veces)} \]
\end{boxsummary}

\begin{boxsummary}
\begin{multicols}{2}
\begin{align*}
a^m \cdot a^n &= a^{m+n} \\[6pt]
\dfrac{a^m}{a^n} &= a^{m-n} \\[6pt]
\left(a^m\right)^n &= a^{m \cdot n}
\end{align*}

\columnbreak

\begin{align*}
(a \cdot b)^n &= a^n \cdot b^n \\[6pt]
\left(\dfrac{a}{b}\right)^n &= \dfrac{a^n}{b^n} \\[6pt]
a^0 &= 1
\end{align*}
\end{multicols}
\end{boxsummary}

\begin{boxsummary}
\begin{multicols}{2}
\begin{align*}
a^{-n} &= \dfrac{1}{a^n} \\[10pt]
\left(\dfrac{a}{b}\right)^{-n} &= \left(\dfrac{b}{a}\right)^n
\end{align*}

\columnbreak

\begin{align*}
a^{-1} &= \dfrac{1}{a} \\[10pt]
\dfrac{1}{a^{-n}} &= a^n
\end{align*}
\end{multicols}
\end{boxsummary}

\begin{boxsummary}
\begin{multicols}{2}
\begin{align*}
a^m \cdot a^n &= a^{m+n} \\[8pt]
\dfrac{a^m}{a^n} &= a^{m-n} \\[8pt]
\left(a^m\right)^n &= a^{m \cdot n} \\[8pt]
(a \cdot b)^n &= a^n \cdot b^n \\[8pt]
\left(\dfrac{a}{b}\right)^n &= \dfrac{a^n}{b^n}
\end{align*}

\columnbreak

\begin{align*}
a^0 &= 1 \\[8pt]
a^{-n} &= \dfrac{1}{a^n} \\[8pt]
\left(\dfrac{a}{b}\right)^{-n} &= \left(\dfrac{b}{a}\right)^n \\[8pt]
a^{-1} &= \dfrac{1}{a} \\[8pt]
\dfrac{1}{a^{-n}} &= a^n
\end{align*}
\end{multicols}
\end{boxsummary}

CONSEJO: Para dividir dos potencias de la misma base entera se recomienda restar el mayor menos el
menor exponente, dejando la potencia donde estaba el mayor exponente (De esta forma
evitamos exponentes negativos)




\begin{enumerate}
    \item Simplificar, utilizando las propiedades de las potencias, dejando el \textbf{resultado como potencia única}, sin usar la calculadora.
        \begin{enumerate}
            \item $\dfrac{14^6}{7^6} =$
            \item $\dfrac{5^6 \cdot 5^7}{5^{11}} =$
            \item $2^2 \cdot \left(2^3\right)^2 =$
            \item $\dfrac{3^8}{\left(3^2\right)^2 \cdot 3} =$
            \item $\left(2^2\right)^4 \cdot a^2 \cdot \left(a^3\right)^2 =$
            \item $\left(2^5 \cdot 7^5\right)^0 =$
            \item $\left(\dfrac{3}{5}\right)^3 \cdot \left(\dfrac{3}{5}\right)^6 =$
            \item $\left(-\dfrac{2}{3}\right) \cdot \left(-\dfrac{2}{3}\right)^9 =$
            \item $\dfrac{(-3)^6 \cdot 3^{-1} \cdot 9^{-2}}{\left[\left(3^2\right)^3\right]^{-2} \cdot 27 \cdot \left(\dfrac{1}{3}\right)^{-2}} =$ \hfill \textit{(Soluc: 6561)}
            \item $\dfrac{3^{-2} \cdot 7^2 \cdot 3 \cdot 7^{-4} \cdot 3^5}{7^3 \cdot 3^{-1} \cdot 7^{-5} \cdot 3^4} =$ \hfill \textit{(Soluc: 3)}
            \item $\dfrac{3^8 \cdot 7^{-1} \cdot 5^2 \cdot 7^3 \cdot 3^{-2}}{7^4 \cdot 5^{-1} \cdot 3^5 \cdot 5^3 \cdot 7^{-2}} =$ \hfill \textit{(Soluc: 3)}
            \item $\dfrac{15^2 \cdot 5^{-2} \cdot 5^3 \cdot 45^2}{\left(5^3\right)^2 \cdot 27 \cdot 3^{-2}} =$ \hfill \textit{(Soluc: 243/5)}
            \item $\dfrac{\left\{\left[(-27)^2\right]^{-3}\right\}^{-2} \cdot 81^{-1}}{3^{-6} \cdot 3^0} =$ \hfill \textit{(Soluc: $3^{38}$)}
            \item $(-2)^3 - \dfrac{2}{3} \cdot (-2)^2 - \dfrac{-2}{4} + 10 =$ \hfill \textit{(Soluc: -1/6)}
            \item $\dfrac{\left[\left(5^2\right)^{-3}\right]^4 \cdot 75^{12}}{1^{-7} \cdot \left(\dfrac{1}{3}\right)^6 \cdot \left(\dfrac{1}{3}\right)^{18} \cdot \left(\dfrac{1}{3}\right)^0} =$ \hfill \textit{(Soluc: 1)}
            \item $\dfrac{9^2 \cdot 2^9 \cdot 8^{-1}}{4^3 \cdot 12 \cdot 3^{-6} \cdot 8} =$ \hfill \textit{(Soluc: $2^7 \cdot 3^9$)}
            \item  $\dfrac{2^3 \cdot 8^{-3} \cdot 12^{-1} \cdot (-3)^2}{6^2 \cdot 16^{-2} \cdot 3^{-3}} =$ \hfill \textit{(Soluc: 9/4)}
            \item $\dfrac{6^4 \cdot 9^2 \cdot 2^{-4} \cdot 3^{-5} \cdot 2^{-1}}{18^3 \cdot 2^{-5} \cdot 3^6 \cdot \left(3^3\right)^{-3}} =$ \hfill \textit{(Soluc: 2)}
            \item $\dfrac{4^4 \cdot 8^{-1} \cdot 16^2}{\left(\dfrac{1}{2}\right)^3 \cdot 8^6} =$ \hfill \textit{(Soluc: 1/4)}
            \item $\dfrac{\dfrac{\left(5^2 \cdot 5^3 \cdot 5^{-4}\right)^2}{\left(5^{-2} \cdot 5^{-3} \cdot 5^4\right)^3}}{\left[\left(\dfrac{1}{5}\right)^2 \cdot \left(\dfrac{1}{5}\right)^4\right]^4} =$ \hfill \textit{(Soluc: 1/125)}
            \item $\dfrac{\left(\dfrac{3}{2}\right)^{-2} \cdot \left(\dfrac{5}{3}\right)^2 \cdot \left(\dfrac{2}{5}\right)^{-1}}{\left(\dfrac{5}{2}\right)^2 \cdot \left(\dfrac{3}{5}\right)^{-3} \cdot 8 \cdot \left(\dfrac{2}{3}\right)^{-1} \cdot 3^{-2}} =$ \hfill \textit{(Soluc: 2/15)}
            \item $\dfrac{\left[\left(\dfrac{2}{3}\right)^2\right]^{-2} \cdot \left(\dfrac{4}{9}\right)^1}{(-2)^6 \cdot \left(2^{-5} \cdot 4\right)^2 \cdot 36} =$ \hfill \textit{(Soluc: 1/81)}
            \item $\dfrac{\left(2^{-4} \cdot 4^3\right)^2 \cdot 5 \cdot 5^0}{100^2 \cdot \left(5^2\right)^{-3}} =$ \hfill \textit{(Soluc: 125)}
            \item $\left(\dfrac{8}{9}\right)^{-2} \cdot \left[\left(-\dfrac{2}{3}\right)^2\right]^3 \cdot \left(\dfrac{3}{2}\right)^{-1} \cdot 9 =$ \hfill \textit{(Soluc: 2/3)}
            \item $\dfrac{\left(\dfrac{3}{2}\right)^0 \cdot \left(\dfrac{2}{3}\right)^{-1}}{\left[\left(\dfrac{3}{2}\right)^{-2}\right]^3 \cdot \left(-\dfrac{3}{2}\right)^4} =$ \hfill \textit{(Soluc: 27/8)}
                    
        \end{enumerate}
     \item Completa las siguientes tables que resumen los casos de cálculo con potencias:
            %\begin{boxsummary}
                \begin{center}
                \renewcommand{\arraystretch}{1.5}
                \begin{tabular}{cc|c|c|}
                \cline{3-4}
                & & \multicolumn{2}{c|}{\textbf{EXPONENTE}} \\ \cline{3-4} 
                & & \textbf{POSITIVO} & \textbf{NEGATIVO} \\ \hline
                \multicolumn{1}{|c|}{\multirow{3}{*}{\rotatebox[origin=c]{90}{\textbf{BASE ENTERA}}}} & \rotatebox[origin=c]{90}{\textbf{POSITIVA}} & $2^3 =$ & $2^{-3} =$ \\ \cline{2-4} 
                \multicolumn{1}{|c|}{} & \multirow{2}{*}{\rotatebox[origin=c]{90}{\textbf{NEGATIVA}}} & $(-2)^3 =$ & $(-2)^{-3} =$ \\
                \multicolumn{1}{|c|}{} & & $(-2)^4 =$ & $(-2)^{-4} =$ \\ \hline
                \end{tabular}
                \hfill
                \begin{tabular}{cc|c|c|}
                \cline{3-4}
                & & \multicolumn{2}{c|}{\textbf{EXPONENTE}} \\ \cline{3-4} 
                & & \textbf{POSITIVO} & \textbf{NEGATIVO} \\ \hline
                \multicolumn{1}{|c|}{\multirow{3}{*}{\rotatebox[origin=c]{90}{\textbf{BASE FRACCIONARIA}}}} & \rotatebox[origin=c]{90}{\textbf{POSITIVA}} & $\left(\dfrac{2}{3}\right)^3 =$ & $\left(\dfrac{2}{3}\right)^{-3} =$ \\ \cline{2-4} 
                \multicolumn{1}{|c|}{} & \multirow{2}{*}{\rotatebox[origin=c]{90}{\textbf{NEGATIVA}}} & $\left(-\dfrac{2}{3}\right)^3 =$ & $\left(-\dfrac{2}{3}\right)^{-3} =$ \\
                \multicolumn{1}{|c|}{} & & $\left(-\dfrac{2}{3}\right)^4 =$ & $\left(-\dfrac{2}{3}\right)^{-4} =$ \\ \hline
                \end{tabular}
                \end{center}

            %\end{boxsummary}
        \item Realiza las siguientes operaciones sin usar la calculadora
        \begin{enumerate}
                \item[\textbf{a)}] $2,5 \cdot 10^7 + 3,6 \cdot 10^7 =$
                \item[\textbf{b)}] $4,6 \cdot 10^{-8} + 5,4 \cdot 10^{-8} =$
                \item[\textbf{c)}] $1,5 \cdot 10^6 + 2,4 \cdot 10^5 =$
                \item[\textbf{d)}] $2,3 \cdot 10^9 + 3,25 \cdot 10^{12} =$
                \item[\textbf{e)}] $3,2 \cdot 10^8 - 1,1 \cdot 10^8 =$
                \item[\textbf{f)}] $4,25 \cdot 10^7 - 2,14 \cdot 10^5 =$
                \item[\textbf{g)}] $7,28 \cdot 10^{-3} - 5,12 \cdot 10^{-3} =$
                \item[\textbf{h)}] $(2 \cdot 10^9) \cdot (3,5 \cdot 10^7) =$
                \item[\textbf{i)}] $\frac{8,4 \cdot 10^9}{2 \cdot 10^7} =$
                \item[\textbf{j)}] $\frac{(3,2 \cdot 10^{-3}) \cdot (4 \cdot 10^5)}{2 \cdot 10^{-8}} =$
                \item[\textbf{k)}] $(2 \cdot 10^5)^2 =$
        \end{enumerate}

\end{enumerate}

\end{document}